%%%%%%%%%%%%%%%%
%%% deut 6 %%%
%%%%%%%%%%%%%%%%

\Chapter{6}

\Verse{1}
{
	\hDtrA{וְזֹאת הַמִּצְוָה הַחֻקִּים וְהַמִּשְׁפָּטִים אֲשֶׁר צִוָּה יְהֹוָה אֱלֹהֵיכֶם לְלַמֵּד אֶתְכֶם לַעֲשׂוֹת בָּאָרֶץ אֲשֶׁר אַתֶּם עֹבְרִים שָׁמָּה לְרִשְׁתָּהּ׃}
}
{
	\eDtrA{
		6 “And this is the commandment, the laws, and the
		judgments that YHWH, your {\God}, commanded to teach you to
		do in the land to which you’re crossing to take possession
		of it,
	}
}
{
}

\Verse{2}
{
	\hDtrA{לְמַעַן תִּירָא אֶת יְהֹוָה אֱלֹהֶיךָ לִשְׁמֹר אֶת כּל חֻקֹּתָיו וּמִצְוֺתָיו אֲשֶׁר אָנֹכִי מְצַוֶּךָ אַתָּה וּבִנְךָ וּבֶן בִּנְךָ כֹּל יְמֵי חַיֶּיךָ וּלְמַעַן יַאֲרִכֻן יָמֶיךָ׃}
}
{
	\eDtrA{
		so that you’ll fear YHWH, your {\God}, to observe all His
		laws and His commandments that I’m commanding you: you and
		your child and your child’s child, all the days of your
		life, and so that your days will be extended.
	}
}
{
}

\Verse{3}
{
	\hDtrA{וְשָׁמַעְתָּ יִשְׂרָאֵל וְשָׁמַרְתָּ לַעֲשׂוֹת אֲשֶׁר יִיטַב לְךָ וַאֲשֶׁר תִּרְבּוּן מְאֹד כַּאֲשֶׁר דִּבֶּר יְהֹוָה אֱלֹהֵי אֲבֹתֶיךָ לָךְ אֶרֶץ זָבַת חָלָב וּדְבָשׁ׃}
}
{
	\eDtrA{
		And you shall listen, Israel, and be watchful to do it,
		that it will be good for you and that you’ll multiply very
		much, as YHWH, your fathers’ {\God}, spoke to you: a land
		flowing with milk and honey.
	}
}
{
}

\Verse{4}
{
	\hDtrA{שְׁמַע יִשְׂרָאֵל יְהֹוָה אֱלֹהֵינוּ יְהֹוָה אֶחָד׃}
}
{
	\eDtrA{“Listen, Israel: YHWH is our {\God}. YHWH is one.}
}
{
%	\footnote{
%		Listen. Forms of this word (Hebrew [image][image]ma‘) occur about eighty
%		times in Deuteronomy. This emphasis on listening goes together with the
%		title and theme of the book: words, Moses’ speech. The great events of
%		creation, patriarchs, exodus, Sinai, and the wilderness are at an end.
%		What remains now is the words: the story of what happened, and the
%		commandments to be fulfilled. Israel’s task now—and in its future
%		generations—is to listen, to learn the story and to do the commandments.
%		And so this passage, probably the most commonly known and frequently
%		repeated words in the Bible by Jews, begins with the commandment that
%		necessarily precedes all others: Listen. Footnote 6:4. YHWH is one. This
%		is a supremely monotheistic statement, especially when taken together
%		with the preceding statements: “He is God” and “There is no other
%		outside of Him” and “There isn’t another.” More evidence that monotheism
%		is essential in Deuteronomy, in its early (preexilic) layers, will
%		follow. Monotheism was an early concept in biblical Israel. Some take
%		these words to mean, “YHWH is our God, YHWH alone.” But “alone” is an
%		extremely uncommon meaning for the Hebrew word here: ’e[image][image]d.
%		Footnote 6:4. YHWH is one. In comparing Israel’s monotheism to pagan
%		religion, we must appreciate that the difference between one and many is
%		not the same sort of thing as the difference between two and three or
%		between six and twenty. It is not numerical. It is a different concept
%		of what a god is. A God who is outside of nature, known through acts in
%		history, a creator, unseeable, without a mate, who makes legal covenants
%		with humans, who is one, is a revolution in religious conception.
%	}
}

\Verse{5}
{
	\hDtrA{וְאָהַבְתָּ אֵת יְהֹוָה אֱלֹהֶיךָ בְּכל לְבָבְךָ וּבְכל נַפְשְׁךָ וּבְכל מְאֹדֶךָ׃}
}
{
	\eDtrA{
		And you shall love YHWH, your {\God}, with all your heart and
		with all your soul and with all your might.
	}
}
{
}

\Verse{6}
{
	\hDtrA{וְהָיוּ הַדְּבָרִים הָאֵלֶּה אֲשֶׁר אָנֹכִי מְצַוְּךָ הַיּוֹם עַל לְבָבֶךָ׃}
}
{
	\eDtrA{
		And these words that I command you today shall be on your
		heart.
	}
}
{
}

\Verse{7}
{
	\hDtrA{וְשִׁנַּנְתָּם לְבָנֶיךָ וְדִבַּרְתָּ בָּם בְּשִׁבְתְּךָ בְּבֵיתֶךָ וּבְלֶכְתְּךָ בַדֶּרֶךְ וּבְשׁכְבְּךָ וּבְקוּמֶךָ׃}
}
{
	\eDtrA{
		And you shall impart them to your children, and you shall
		speak about them when you sit in your house and when you
		go in the road and when you lie down and when you get up.
	}
}
{
}

\Verse{8}
{
	\hDtrA{וּקְשַׁרְתָּם לְאוֹת עַל יָדֶךָ וְהָיוּ לְטֹטָפֹת בֵּין עֵינֶיךָ׃}
}
{
	\eDtrA{
		And you shall bind them for a sign on your hand, and they
		shall become bands between your eyes.
	}
}
{
%	\footnote{
%		bind them on your hand. This came to be taken literally, requiring one
%		to wear boxes (tephillin) on one’s arm and head containing passages from
%		the Torah. In the Tanak, however, this expression is meant figuratively,
%		meaning to keep these teachings at hand. Thus, elsewhere it refers to
%		binding religious teachings to one’s neck (Prov 3:3), one’s heart
%		(6:21), and one’s fingers (7:3). Footnote 6:8. bands between your eyes.
%		This, too, is meant figuratively, meaning to keep these teachings right
%		before one’s eyes. Elsewhere it is used to refer to the law of
%		redemption of the firstborn (Exod 13:16), where there is no connection
%		to tephillin.
%	}
}

\Verse{9}
{
	\hDtrA{וּכְתַבְתָּם עַל מְזֻזוֹת בֵּיתֶךָ וּבִשְׁעָרֶיךָ׃}
}
{
	\eDtrA{
		And you shall write them on the doorposts of your house
		and in your gates.
	}
}
{
}

\Verse{10}
{
	\hDtrA{וְהָיָה כִּי יְבִיאֲךָ יְהֹוָה אֱלֹהֶיךָ אֶל הָאָרֶץ אֲשֶׁר נִשְׁבַּע לַאֲבֹתֶיךָ לְאַבְרָהָם לְיִצְחָק וּלְיַעֲקֹב לָתֶת לָךְ עָרִים גְּדֹלֹת וְטֹבֹת אֲשֶׁר לֹא בָנִיתָ׃}
}
{
	\eDtrA{
		“And it will be when YHWH, your {\God}, will bring you to the
		land that He swore to your fathers, to Abraham, to Isaac,
		and to Jacob, to give you big and good cities that you
		didn’t build,
	}
}
{
}

\Verse{11}
{
	\hDtrA{וּבָתִּים מְלֵאִים כּל טוּב אֲשֶׁר לֹא מִלֵּאתָ וּבֹרֹת חֲצוּבִים אֲשֶׁר לֹא חָצַבְתָּ כְּרָמִים וְזֵיתִים אֲשֶׁר לֹא נָטָעְתָּ וְאָכַלְתָּ וְשָׂבָעְתָּ׃}
}
{
	\eDtrA{
		and houses filled with everything good that you didn’t
		fill, and cisterns hewed that you didn’t hew, vineyards
		and olives that you did’t plant, and you’ll eat and be
		satisfied,
	}
}
{
}

\Verse{12}
{
	\hDtrA{הִשָּׁמֶר לְךָ פֶּן תִּשְׁכַּח אֶת יְהֹוָה אֲשֶׁר הוֹצִיאֲךָ מֵאֶרֶץ מִצְרַיִם מִבֵּית עֲבָדִים׃}
}
{
	\eDtrA{
		watch yourself in case you’ll forget YHWH, who brought you
		out from the land of Egypt, from a house of slaves.
	}
}
{
}

\Verse{13}
{
	\hDtrA{אֶת יְהֹוָה אֱלֹהֶיךָ תִּירָא וְאֹתוֹ תַעֲבֹד וּבִשְׁמוֹ תִּשָּׁבֵעַ׃}
}
{
	\eDtrA{
		It’s YHWH, your {\God}, whom you’ll fear, and it’s He whom
		you’ll serve, and it’s in His name that you’ll swear.
	}
}
{
}

\Verse{14}
{
	\hDtrA{לֹא תֵלְכוּן אַחֲרֵי אֱלֹהִים אֲחֵרִים מֵאֱלֹהֵי הָעַמִּים אֲשֶׁר סְבִיבוֹתֵיכֶם׃}
}
{
	\eDtrA{
		You shall not go after other gods, from the gods of the
		peoples who are around you,
	}
}
{
}

\Verse{15}
{
	\hDtrA{כִּי אֵל קַנָּא יְהֹוָה אֱלֹהֶיךָ בְּקִרְבֶּךָ פֶּן יֶחֱרֶה אַף יְהֹוָה אֱלֹהֶיךָ בָּךְ וְהִשְׁמִידְךָ מֵעַל פְּנֵי הָאֲדָמָה׃}
}
{
	\eDtrA{
		because YHWH, your {\God}, is a jealous {\God} among you, in
		case the anger of YHWH, your {\God}, will flare at you, and
		He’ll destroy you from the face of the earth.
	}
}
{
}

\Verse{16}
{
	\hDtrA{לֹא תְנַסּוּ אֶת יְהֹוָה אֱלֹהֵיכֶם כַּאֲשֶׁר נִסִּיתֶם בַּמַּסָּה׃}
}
{
	\eDtrA{
		You shall not test YHWH, your {\God}, as you tested at
		Massah.
	}
}
{
}

\Verse{17}
{
	\hDtrA{שָׁמוֹר תִּשְׁמְרוּן אֶת מִצְוֺת יְהֹוָה אֱלֹהֵיכֶם וְעֵדֹתָיו וְחֻקָּיו אֲשֶׁר צִוָּךְ׃}
}
{
	\eDtrA{
		You shall observe the commandments of YHWH, your {\God}, and
		His testimonies and His laws that He commanded you.
	}
}
{
}

\Verse{18}
{
	\hDtrA{וְעָשִׂיתָ הַיָּשָׁר וְהַטּוֹב בְּעֵינֵי יְהֹוָה לְמַעַן יִיטַב לָךְ וּבָאתָ וְיָרַשְׁתָּ אֶת הָאָרֶץ הַטֹּבָה אֲשֶׁר נִשְׁבַּע יְהֹוָה לַאֲבֹתֶיךָ׃}
}
{
	\eDtrA{
		And you shall do what is right and good in YHWH’s eyes so
		it will be good for you, and you’ll come and take
		possession of the good land that YHWH swore to your
		fathers,
	}
}
{
}

\Verse{19}
{
	\hDtrA{לַהֲדֹף אֶת כּל אֹיְבֶיךָ מִפָּנֶיךָ כַּאֲשֶׁר דִּבֶּר יְהֹוָה׃}
}
{
	\eDtrA{
		to push all your enemies from in front of you, as YHWH has
		spoken.
	}
}
{
}

\Verse{20}
{
	\hDtrA{כִּי יִשְׁאָלְךָ בִנְךָ מָחָר לֵאמֹר מָה הָעֵדֹת וְהַחֻקִּים וְהַמִּשְׁפָּטִים אֲשֶׁר צִוָּה יְהֹוָה אֱלֹהֵינוּ אֶתְכֶם׃}
}
{
	\eDtrA{
		“When your child will ask you tomorrow, saying, ‘What are
		the testimonies and the laws and the judgments that YHWH,
		our {\God}, commanded you?’
	}
}
{
}

\Verse{21}
{
	\hDtrA{וְאָמַרְתָּ לְבִנְךָ עֲבָדִים הָיִינוּ לְפַרְעֹה בְּמִצְרָיִם וַיֹּצִיאֵנוּ יְהֹוָה מִמִּצְרַיִם בְּיָד חֲזָקָה׃}
}
{
	\eDtrA{
		then you shall say to your child, ‘We were slaves to
		Pharaoh in Egypt, and YHWH brought us out from Egypt with
		a strong hand.
	}
}
{
}

\Verse{22}
{
	\hDtrA{וַיִּתֵּן יְהֹוָה אוֹתֹת וּמֹפְתִים גְּדֹלִים וְרָעִים בְּמִצְרַיִם בְּפַרְעֹה וּבְכל בֵּיתוֹ לְעֵינֵינוּ׃}
}
{
	\eDtrA{
		And YHWH put great and harsh signs and wonders in Egypt at
		Pharaoh and at all of his household before our eyes.
	}
}
{
}

\Verse{23}
{
	\hDtrA{וְאוֹתָנוּ הוֹצִיא מִשָּׁם לְמַעַן הָבִיא אֹתָנוּ לָתֶת לָנוּ אֶת הָאָרֶץ אֲשֶׁר נִשְׁבַּע לַאֲבֹתֵינוּ׃}
}
{
	\eDtrA{
		And He brought us out from there in order to bring us to
		give us the land that He swore to our fathers.
	}
}
{
}

\Verse{24}
{
	\hDtrA{וַיְצַוֵּנוּ יְהֹוָה לַעֲשׂוֹת אֶת כּל הַחֻקִּים הָאֵלֶּה לְיִרְאָה אֶת יְהֹוָה אֱלֹהֵינוּ לְטוֹב לָנוּ כּל הַיָּמִים לְחַיֹּתֵנוּ כְּהַיּוֹם הַזֶּה׃}
}
{
	\eDtrA{
		And YHWH commanded us to do all these laws, to fear YHWH,
		our {\God}, for our good every day, to keep us alive as we
		are this day.
	}
}
{
}

\Verse{25}
{
	\hDtrA{וּצְדָקָה תִּהְיֶה לָּנוּ כִּי נִשְׁמֹר לַעֲשׂוֹת אֶת כּל הַמִּצְוָה הַזֹּאת לִפְנֵי יְהֹוָה אֱלֹהֵינוּ כַּאֲשֶׁר צִוָּנוּ׃}
}
{
	\eDtrA{
		And we’ll have virtue when we are watchful to do all of
		this commandment in front of YHWH, our {\God}, as He
		commanded us.’
	}
}
{
%	\footnote{
%		we’ll have virtue. Again an issue from Genesis returns in Moses’
%		farewell in Deuteronomy. The initial covenant with Abraham is made on
%		the premise of Abraham’s virtue ([image][image]d[image]q[image]h, Gen
%		15:6), and God later says of Abraham that “I’ve known him for the
%		purpose that he’ll command his children and his house after him, and
%		they’ll observe YHWH’s way, to do virtue … ” (18:19). His grandson Jacob
%		later says that “my virtue will answer for me” (30:33). And then the
%		word does not occur again in the Torah until this speech in Deuteronomy,
%		where Moses repeatedly tells the people how they can come to have
%		virtue.
%	}
}
